% GENERATED FILE. Edit canonical lessons, then run npm run generate:books.
% canonical-source-hash: fnv1a64:563d05ce5b3f391c
% canonical-lessons: BN-C346-prarthi, BN-C346-rashtropoti, BN-C346-bhotar, BN-C346-podotyag-kora, BN-C346-shopoth-neoya

\chapter{Standing for Office}
\label{ch:bn-standing-for-office}
\hlchaptermodality{346}

\begin{chapteropening}
\textbf{By the end of this chapter:} \emph{I can talk about who stands for office, who votes and who steps down.}

Complete the last lesson of chapter 346: I can talk about who stands for office, who votes and who steps down.
\end{chapteropening}

\section[\texorpdfstring{\bn{প্রার্থী} (prārthī)}{প্রার্থী (prārthī)}]{\bn{প্রার্থী} (prārthī) — a candidate}

\label{lesson:BN-C346-prarthi}



\begin{quote}
Before the new one: say the Bengali for to approve, then the Bengali for to repeal.
\end{quote}

\subsection*{You'll want to know: \bn{প্রার্থী}}
\textbf{\bn{প্রার্থী}} — \emph{prārthī} — \textquotedblleft{}a candidate\textquotedblright{}.

\begin{quote}
\bn{এই} \bn{নির্বাচনে} \bn{চারজন} \bn{প্রার্থী} \bn{আছেন।}
\end{quote}

\begin{quote}
\emph{ei nirbāchône chārjôn prārthī āchhen.}
\end{quote}

\begin{quote}
There are four candidates in this election.
\end{quote}

\begin{cousinweb}
\textbf{\bn{প্রার্থী}} (\emph{prārthī}) is Sanskrit \emph{prārthin}, one who wishes for something (Monier-Williams). Bengali uses it for one who seeks an office or a job, a candidate.
\end{cousinweb}

\begin{grammarlens}[title={What you've built}]
One more word for giving reasons and opinions.
\end{grammarlens}

\subsection*{Your turn}
\begin{itemize}
\raggedright
  \item \emph{Say it:} \emph{prārthī}
  \item \emph{Say it:} \emph{prārthī}, once more
  \item \emph{Say it:} \emph{ei nirbāchône chārjôn prārthī āchhen.} — There are four candidates in this election.
  \item \emph{Recall:} say the Bengali for to approve, then the Bengali for to repeal, then say \emph{prārthī} again
\end{itemize}

\begin{tcolorbox}[breakable,colback=teal!4,colframe=teal!35!black,title={Before you move on}]
What does \bn{প্রার্থী} mean? (\textquotedblleft{}A candidate\textquotedblright{}.) Say it once more.
\end{tcolorbox}
\section[\texorpdfstring{\bn{রাষ্ট্রপতি} (rāshṭrôpôti)}{রাষ্ট্রপতি (rāshṭrôpôti)}]{\bn{রাষ্ট্রপতি} (rāshṭrôpôti) — a president (of a country)}

\label{lesson:BN-C346-rashtropoti}



\begin{quote}
Before the new one: say the Bengali for to repeal, then the Bengali for a candidate.
\end{quote}

\subsection*{You'll want to know: \bn{রাষ্ট্রপতি}}
\textbf{\bn{রাষ্ট্রপতি}} — \emph{rāshṭrôpôti} — \textquotedblleft{}a president (of a country)\textquotedblright{}.

\begin{quote}
\bn{রাষ্ট্রপতি} \bn{সংসদে} \bn{বক্তব্য} \bn{রাখলেন।}
\end{quote}

\begin{quote}
\emph{rāshṭrôpôti sôngsôde bôktôbyô rākhlen.}
\end{quote}

\begin{quote}
The president addressed parliament.
\end{quote}

\begin{cousinweb}
\textbf{\bn{রাষ্ট্রপতি}} (\emph{rāshṭrôpôti}) is Sanskrit \emph{rāṣṭrapati}, lord of a realm (Monier-Williams). The modern Bengali sense, an elected head of state, is a twentieth-century use.
\end{cousinweb}

\begin{grammarlens}[title={What you've built}]
One more word for giving reasons and opinions.
\end{grammarlens}

\subsection*{Your turn}
\begin{itemize}
\raggedright
  \item \emph{Say it:} \emph{rāshṭrôpôti}
  \item \emph{Say it:} \emph{rāshṭrôpôti}, once more
  \item \emph{Say it:} \emph{rāshṭrôpôti sôngsôde bôktôbyô rākhlen.} — The president addressed parliament.
  \item \emph{Recall:} say the Bengali for to repeal, then the Bengali for a candidate, then say \emph{rāshṭrôpôti} again
\end{itemize}

\begin{tcolorbox}[breakable,colback=teal!4,colframe=teal!35!black,title={Before you move on}]
What does \bn{রাষ্ট্রপতি} mean? (\textquotedblleft{}A president (of a country)\textquotedblright{}.) Say it once more.
\end{tcolorbox}
\section[\texorpdfstring{\bn{ভোটার} (bhoṭār)}{ভোটার (bhoṭār)}]{\bn{ভোটার} (bhoṭār) — a voter}

\label{lesson:BN-C346-bhotar}



\begin{quote}
Before the new one: say the Bengali for a candidate, then the Bengali for a president.
\end{quote}

\subsection*{You'll want to know: \bn{ভোটার}}
\textbf{\bn{ভোটার}} — \emph{bhoṭār} — \textquotedblleft{}a voter\textquotedblright{}.

\begin{quote}
\bn{এই} \bn{গ্রামে} \bn{দুই} \bn{হাজার} \bn{ভোটার} \bn{আছে।}
\end{quote}

\begin{quote}
\emph{ei grāme dui hājār bhoṭār āchhe.}
\end{quote}

\begin{quote}
There are two thousand voters in this village.
\end{quote}

\begin{cousinweb}
\textbf{\bn{ভোটার}} (\emph{bhoṭār}) is English \emph{voter}, beside \bn{ভোট} from English \emph{vote}; Bengali has no v, so it is heard as bh.
\end{cousinweb}

\begin{grammarlens}[title={What you've built}]
One more word for giving reasons and opinions.
\end{grammarlens}

\subsection*{Your turn}
\begin{itemize}
\raggedright
  \item \emph{Say it:} \emph{bhoṭār}
  \item \emph{Say it:} \emph{bhoṭār}, once more
  \item \emph{Say it:} \emph{ei grāme dui hājār bhoṭār āchhe.} — There are two thousand voters in this village.
  \item \emph{Recall:} say the Bengali for a candidate, then the Bengali for a president, then say \emph{bhoṭār} again
\end{itemize}

\begin{tcolorbox}[breakable,colback=teal!4,colframe=teal!35!black,title={Before you move on}]
What does \bn{ভোটার} mean? (\textquotedblleft{}A voter\textquotedblright{}.) Say it once more.
\end{tcolorbox}
\section[\texorpdfstring{\bn{পদত্যাগ} \bn{করা} (pôdôtyāg kôrā)}{পদত্যাগ করা (pôdôtyāg kôrā)}]{\bn{পদত্যাগ} \bn{করা} (pôdôtyāg kôrā) — to resign}

\label{lesson:BN-C346-podotyag-kora}



\begin{quote}
Before the new one: say the Bengali for a president, then the Bengali for a voter.
\end{quote}

\subsection*{You'll want to know: \bn{পদত্যাগ} \bn{করা}}
\textbf{\bn{পদত্যাগ} \bn{করা}} — \emph{pôdôtyāg kôrā} — \textquotedblleft{}to resign\textquotedblright{}.

\begin{quote}
\bn{দুর্নীতির} \bn{অভিযোগের} \bn{পরে} \bn{মন্ত্রী} \bn{পদত্যাগ} \bn{করলেন।}
\end{quote}

\begin{quote}
\emph{durnītir ôbhijoger pôre môntrī pôdôtyāg kôrlen.}
\end{quote}

\begin{quote}
The minister resigned after the corruption charges.
\end{quote}

Past (he/she …): \bn{তিনি} \bn{পদত্যাগ} \bn{করলেন}, \emph{tini pôdôtyāg kôrlen}.

\begin{cousinweb}
\textbf{\bn{পদত্যাগ} \bn{করা}} (\emph{pôdôtyāg kôrā}) is Sanskrit \emph{pada}, a position, with \emph{tyāga}, leaving or giving up (Monier-Williams), and \bn{করা}: to give up one's post.
\end{cousinweb}

\begin{grammarlens}[title={What you've built}]
One more word for giving reasons and opinions.
\end{grammarlens}

\subsection*{Your turn}
\begin{itemize}
\raggedright
  \item \emph{Say it:} \emph{pôdôtyāg kôrā}
  \item \emph{Say it:} \emph{pôdôtyāg kôrā}, once more
  \item \emph{Say it:} \emph{durnītir ôbhijoger pôre môntrī pôdôtyāg kôrlen.} — The minister resigned after the corruption charges.
  \item \emph{Recall:} say the Bengali for a president, then the Bengali for a voter, then say \emph{pôdôtyāg kôrā} again
\end{itemize}

\begin{tcolorbox}[breakable,colback=teal!4,colframe=teal!35!black,title={Before you move on}]
What does \bn{পদত্যাগ} \bn{করা} mean? (\textquotedblleft{}To resign\textquotedblright{}.) Say it once more.
\end{tcolorbox}
\section[\texorpdfstring{\bn{শপথ} \bn{নেওয়া} (shôpôth neoyā)}{শপথ নেওয়া (shôpôth neoyā)}]{\bn{শপথ} \bn{নেওয়া} (shôpôth neoyā) — to take an oath; to be sworn in}

\label{lesson:BN-C346-shopoth-neoya}



\begin{quote}
Before the new one: say the Bengali for a voter, then the Bengali for to resign.
\end{quote}

\subsection*{You'll want to know: \bn{শপথ} \bn{নেওয়া}}
\textbf{\bn{শপথ} \bn{নেওয়া}} — \emph{shôpôth neoyā} — \textquotedblleft{}to take an oath; to be sworn in\textquotedblright{}.

\begin{quote}
\bn{নতুন} \bn{রাষ্ট্রপতি} \bn{আজ} \bn{শপথ} \bn{নিলেন।}
\end{quote}

\begin{quote}
\emph{nôtun rāshṭrôpôti āj shôpôth nilen.}
\end{quote}

\begin{quote}
The new president was sworn in today.
\end{quote}

Past (he/she …): \bn{তিনি} \bn{শপথ} \bn{নিলেন}, \emph{tini shôpôth nilen}.

\begin{cousinweb}
\textbf{\bn{শপথ} \bn{নেওয়া}} (\emph{shôpôth neoyā}) is Sanskrit \emph{śapatha}, an oath (Monier-Williams), with \bn{নেওয়া} 'to take': to take an oath, as a new office-holder does.
\end{cousinweb}

\begin{grammarlens}[title={What you've built}]
One more word for giving reasons and opinions.
\end{grammarlens}

\subsection*{Your turn}
\begin{itemize}
\raggedright
  \item \emph{Say it:} \emph{shôpôth neoyā}
  \item \emph{Say it:} \emph{shôpôth neoyā}, once more
  \item \emph{Say it:} \emph{nôtun rāshṭrôpôti āj shôpôth nilen.} — The new president was sworn in today.
  \item \emph{Recall:} say the Bengali for a voter, then the Bengali for to resign, then say \emph{shôpôth neoyā} again
\end{itemize}

\begin{tcolorbox}[breakable,colback=teal!4,colframe=teal!35!black,title={Before you move on}]
What does \bn{শপথ} \bn{নেওয়া} mean? (\textquotedblleft{}To take an oath; to be sworn in\textquotedblright{}.) Say it once more.
\end{tcolorbox}
